% Capítulo 4: Desarrollo del Proyecto

\section{Análisis del sistema}

\subsection{Formalización matemática de los caminos}

Sea una red neuronal feed-forward con activación ReLU. Para una entrada $x \in \mathbb{R}^n$, la
activación de cada neurona en cada capa es no-negativa (por definición de ReLU). Se dice que la
neurona $j$ de la capa $l$ está \textbf{activa} para $x$ si su pre-activación es estrictamente
positiva, e \textbf{inactiva} si es cero.

Un \textbf{camino} (\textit{path}) $p = (j_0, j_1, \ldots, j_L)$ es una secuencia de índices de
neuronas, una por capa, tal que la neurona $j_l$ de la capa $l$ está conectada con la neurona
$j_{l+1}$ de la capa $l+1$. El \textbf{valor del coeficiente} de un camino para una entrada $x$ es
el producto de:
\begin{itemize}
  \item La activación de la neurona de entrada del camino (neurona $j_0$), y
  \item El indicador de actividad de cada neurona intermedia (1 si activa, 0 si no).
\end{itemize}

Dado el conjunto de todos los caminos posibles entre la capa \texttt{initial\_layer} y la salida, y
dado un conjunto de muestras $\{(x_i, y_i)\}$, el sistema lineal $Ax = b$ se construye de la
siguiente forma:
\begin{itemize}
  \item Cada fila de $A$ corresponde a una muestra $x_i$ y contiene los coeficientes de los
        caminos activos para esa muestra.
  \item $b$ contiene las salidas deseadas $y_i$.
  \item $x$ son los pesos de los caminos (incógnitas del problema).
\end{itemize}

La solución $x^*$ del sistema proporciona los pesos óptimos que minimizan el error cuadrático entre
la salida de la red reconstruida y las salidas deseadas.

\subsection{Casos de uso principales}

\paragraph{CU-01: Entrenar Qsimov (primera ronda)}
\begin{itemize}
  \item Actor: usuario de la librería
  \item Precondición: modelo Keras/PyTorch pre-entrenado disponible
  \item Flujo: (1) crear \texttt{PathSelector} con \texttt{initial\_layer}, (2) crear
        \texttt{QsimovLinearSystem}, (3) llamar a \texttt{fit(X\_train, Y\_train)}
  \item Postcondición: \texttt{solutions\_} contiene los pesos de los caminos resolventes
\end{itemize}

\paragraph{CU-02: Actualizar el modelo sin olvido (ronda N)}
\begin{itemize}
  \item Actor: usuario
  \item Precondición: \texttt{QsimovLinearSystem} entrenado en rondas anteriores
  \item Flujo: llamar a \texttt{fit(X\_nueva, Y\_nueva)}, \textit{sin} llamar
        previamente a \texttt{reset\_equations()}
  \item Postcondición: \texttt{solutions\_} incorpora todas las muestras de todas las rondas
\end{itemize}

\paragraph{CU-03: Predecir con el modelo actualizado}
\begin{itemize}
  \item Actor: usuario
  \item Precondición: \texttt{QsimovLinearSystem} con \texttt{solutions\_} calculadas
  \item Flujo: llamar a \texttt{predict(X\_test)}
  \item Postcondición: se devuelven las predicciones del modelo Qsimov
\end{itemize}

\paragraph{CU-04: Guardar y recuperar el modelo}
\begin{itemize}
  \item Actor: usuario
  \item Guardar: \texttt{qls.save("modelo.qsi")} \\
        Recuperar: \texttt{qls = KerasQsimov\-Linear\-System.load("modelo.qsi")}
\end{itemize}

\subsection{Modelado de datos}

Los datos que gestiona el sistema son:
\begin{itemize}
  \item \texttt{all\_paths\_}: lista de arrays numpy de forma \texttt{(n\_paths\_layer, path\_length)},
        los índices que forman cada camino entre capas.
  \item \texttt{output\_masks\_}: array booleano
        \texttt{(n\_outputs, n\_paths\_total)};
        indica qué caminos llegan a cada neurona de salida.
  \item \texttt{equations\_}: lista de sistemas \texttt{[A, b]} por output, donde \texttt{A} tiene
        forma \texttt{(n\_equations, n\_paths\_output)}.
  \item \texttt{solutions\_}: lista de arrays numpy de forma \texttt{(n\_paths\_output,)},
        pesos de los caminos resueltos.
\end{itemize}

\section{Diseño del sistema}

\subsection{Diseño de interfaces de usuario}

La API pública de la librería está diseñada para ser coherente con las convenciones de scikit-learn
(\texttt{fit} / \texttt{predict}) y extensible a nuevos frameworks. La interfaz mínima para el
usuario es:

\begin{lstlisting}[language=Python, caption={API mínima -- variante Keras.}]
# Keras
from qsimov.keras_path_selector import KerasPathSelector
from qsimov.keras_qsimov_linear_system import (
    KerasQsimovLinearSystem)

path_selector = KerasPathSelector(model, initial_layer=-1)
qls = KerasQsimovLinearSystem(
    path_selector, solver="back_substitution")
qls.fit(x_train, y_train)
y_pred = qls.predict(x_test)
\end{lstlisting}

\begin{lstlisting}[language=Python, caption={API mínima -- variante PyTorch.}]
# PyTorch
from qsimov.pytorch_path_selector import PytorchPathSelector
from qsimov.pytorch_qsimov_gradient import PytorchQsimovGradient

path_selector = PytorchPathSelector(
    model, input_shape=(1, 28, 28), initial_layer=-4)
qg = PytorchQsimovGradient(path_selector)
qg.fit(x_train, y_train, epochs=5,
       optimizer=lambda p: torch.optim.Adam(p, lr=1e-3))
\end{lstlisting}

\subsection{Diseño de procesos: extracción de caminos}

El proceso de extracción de caminos para una capa densa con matriz de pesos
$W \in \mathbb{R}^{m \times n}$ funciona de la siguiente forma:

\begin{enumerate}
  \item Para cada neurona de entrada $i$ (incluyendo el bias como neurona $0$): se recorren todas
        las neuronas de salida $j$ con $W[i, j] \neq 0$. Cada par $(i, j)$ constituye un camino
        elemental de longitud 2.
  \item Los caminos de capas sucesivas se combinan mediante \path{combine_paths_left_right}: se
        hace un join de los caminos izquierdos y derechos donde el último índice del camino
        izquierdo coincide con el primero del camino derecho. Los caminos de bias (con primer índice
        0) siempre se incluyen.
  \item El proceso se repite para todas las capas del \texttt{right\_model\_}, acumulando los
        caminos mediante combinación.
\end{enumerate}

Para capas convolucionales, cada posición del mapa de características y cada canal de salida genera
sus propios caminos, aplicando las transformaciones de stride y padding correspondientes.

\subsection{Diseño de la base de datos (persistencia)}

El formato de persistencia \texttt{.qsi} es un directorio con la siguiente estructura:

\begin{lstlisting}[caption={Estructura del directorio de persistencia \texttt{.qsi}.}]
modelo.qsi/
|-- py_objects.pkl           <- atributos escalares y listas (pickle)
|-- numpy_variables.npz      <- arrays numpy comprimidos (np.savez_compressed)
|-- left_model.h5 / .pt      <- left_model_ del PathSelector
|-- right_model.h5 / .pt     <- right_model_ del PathSelector
`-- path_selector.qsi/       <- PathSelector anidado
    |-- py_objects.pkl
    `-- numpy_variables.npz
\end{lstlisting}

Esta estructura separa los arrays numpy de gran tamaño (como \texttt{all\_paths\_} o
\texttt{equations\_}) del resto de atributos de Python, permitiendo una carga eficiente sin
deserializar objetos que no sean necesarios.

\section{Arquitectura interna de Qsimov}

\subsection{Módulo \texttt{paths/}: Generación de caminos (Python + Cython)}

Las operaciones de mayor coste computacional en la generación de caminos están implementadas en
Cython (\texttt{c\_paths.pyx}, \texttt{c\_combine.pyx}, \texttt{c\_dense.pyx},
\texttt{c\_conv.pyx}, \texttt{c\_maxpooling.pyx}). Cython se eligió sobre alternativas como NumPy
vectorizado o C puro porque las operaciones implican bucles anidados sobre arrays de longitud
variable con lógica de control compleja, donde la vectorización NumPy no es directamente aplicable.

La función central \texttt{c\_non\_zero\_input\_select\_paths} filtra, para cada muestra, los
caminos cuya neurona de entrada tiene activación no-nula. Su signatura en Cython es:

\begin{lstlisting}[language=C, caption={Signatura de \texttt{c\_non\_zero\_input\_select\_paths}.}]
def c_non_zero_input_select_paths(
    float[:] flat_inputs_with_bias,   # activaciones de entrada, con bias prepend
    long[:, :] all_paths,             # todos los caminos posibles
) -> tuple[np.ndarray, np.ndarray]:  # (selected_paths, select_mask)
\end{lstlisting}

\subsection{Módulo \texttt{linalg.py}: Resolución del sistema}

El sistema lineal se resuelve siguiendo estos pasos:

\begin{enumerate}
  \item \textbf{Transformación QR} (\texttt{r\_transform}): Se aplica la factorización QR con
        \texttt{numpy.linalg.qr(mode="r")} para obtener la matriz triangular superior $R$. Esto
        reduce el sistema acumulado, que puede tener miles de ecuaciones, a un sistema cuadrado
        de tamaño $n_{paths} \times n_{paths}$, manteniendo solo la información relevante.
  \item \textbf{Back-substitution} (\texttt{c\_back\_substitution}): Se resuelve
        $R \cdot x = Q^T \cdot b$ hacia atrás. Se aplican cutoffs absoluto y relativo para tratar
        los coeficientes casi-nulos como cero exacto, lo que mejora la estabilidad numérica.
  \item \textbf{Shrinkage incremental}: El parámetro \texttt{qr\_shrinkage\_factor} controla cada
        cuántas filas nuevas se aplica la transformación QR para comprimir el sistema en memoria.
\end{enumerate}

La variante de mínimos cuadrados (\texttt{solver="lstsq"}) delega directamente en
\texttt{numpy.linalg.lstsq}, útil cuando el sistema está sub-determinado (más incógnitas que
ecuaciones).

\subsection{Diferencias de implementación entre Keras y PyTorch}

Las dos implementaciones comparten las clases base abstractas \texttt{PathSelector} y
\texttt{QsimovLinearSystem}, pero difieren en los siguientes puntos:

\begin{table}[htbp]
  \centering
  \caption{Diferencias entre las variantes Keras y PyTorch de Qsimov.}
  \label{tab:tabla1}
  \begin{tabular}{p{5cm}p{4.5cm}p{4.5cm}}
    \toprule
    \textbf{Aspecto} & \textbf{Keras} & \textbf{PyTorch} \\
    \midrule
    Inferencia del \texttt{input\_shape}  & Automática (el modelo lo infiere) & Manual (parámetro requerido) \\
    Formato de pesos Dense                & $W \in \mathbb{R}^{in \times out}$ & $W \in \mathbb{R}^{out \times in}$ (se transpone) \\
    Formato de datos conv                 & \texttt{channels\_last}            & \texttt{channels\_first} \\
    Iterador de batches                   & \path{as_tensorflow_dataset()}  & \path{as_pytorch_dataloader()} \\
    Dispositivo del right model           & Siempre CPU (por diseño)           & Siempre CPU (por diseño) \\
    Guardado del modelo                   & Formato \texttt{.h5}               & Formato \texttt{.pt} \\
    Indexación de \texttt{initial\_layer} & Solo capas Dense/Conv (ReLU y Dropout no cuentan) & Todas las capas, incluidas ReLU y Dropout \\
    \bottomrule
  \end{tabular}
\end{table}
